\chapter{Phương pháp}
\label{chap:method}

Chương này phát biểu hình thức bài toán (mục~\ref{sec:m-problem}), định nghĩa thói quen và thuật
toán tìm thói quen (mục~\ref{sec:m-habit}), rồi trình bày lần lượt các khối của quy trình
(Hình~\ref{fig:pipeline}): tiền xử lý, cấu trúc nhà rút ra từ dữ liệu, nhận dạng (người, hoạt động),
các baseline, cách đánh giá thói quen, phân tích lỗi lan truyền và phát hiện bỏ thuốc.

\section{Phát biểu bài toán}
\label{sec:m-problem}

\textbf{Dữ liệu cảm biến.} Một nhà có tập cảm biến $\mathcal S$. Nhật ký là chuỗi sự kiện
$E = (e_1, e_2, \dots)$ sắp theo thời gian, $e_i = (t_i, s_i, v_i)$ với thời điểm $t_i$, cảm biến
$s_i \in \mathcal S$ và giá trị $v_i$. Không có thông tin ai gây ra $e_i$.

\textbf{Người và hoạt động.} Nhà có tập người $\mathcal P$ (ở đây $\{R1, R2\}$). Tập hoạt động
$\mathcal A$ gồm các nhãn của bộ dữ liệu; mỗi hoạt động $a$ gắn với một người
$\rho(a) \in \mathcal P \cup \{H\}$, trong đó $H$ nghĩa là ``chung của hộ'' (nhãn không ghi người,
ví dụ \texttt{Breakfast}). Quy ước của CASAS: nhãn có tiền tố \texttt{R1\_}/\texttt{R2\_} gắn với
người đó.

\textbf{Lần thực hiện.} Một lần thực hiện là bộ $o = (a, t^{s}, t^{e})$: hoạt động $a$ diễn ra
trong $[t^s, t^e]$. Nhãn thật cho tập lần thực hiện $O^\star$; một bộ nhận dạng sinh ra tập dự đoán
$\hat O$.

\textbf{Ngày quan sát.} $\Dd$ là tập ngày có dữ liệu. Với loại ngày
$\tau \in \{\text{tất cả}, \text{ngày thường}, \text{cuối tuần}\}$, $\Dd_\tau \subseteq \Dd$ là các
ngày thuộc loại đó.

\textbf{Bài toán.} Cho $E$ trên các ngày kiểm tra (không có nhãn) và nhãn $O^\star$ trên các ngày huấn
luyện, tìm tập thói quen $\hat\Hb$ của từng người trên các ngày kiểm tra sao cho $\hat\Hb$ khớp với tập
thói quen đáp án $\Hb^\star$ --- thói quen thu được khi áp định nghĩa ở mục~\ref{sec:m-habit} lên
nhãn thật của chính các ngày kiểm tra.

\section{Định nghĩa thói quen}
\label{sec:m-habit}

\begin{definition}[Thói quen]
\label{def:habit}
Cho hoạt động $a$, loại ngày $\tau$, nửa bề rộng cửa sổ $w$ (mặc định 30 phút) và ngưỡng tần suất
$\phi$ (mặc định $0{,}5$). Gọi $N(c)$ là số \emph{ngày phân biệt} trong $\Dd_\tau$ có ít nhất một lần
thực hiện $a$ bắt đầu trong cửa sổ $[c-w, c+w]$ trên vòng 24 giờ. Một \textbf{thói quen} là một tâm
$c$ sao cho
\begin{enumerate}
  \item $N(c) \ge \phi\,|\Dd_\tau|$ (lặp lại ở ít nhất một nửa số ngày);
  \item $c$ là điểm đạt $N$ lớn nhất trong đoạn liên tục các tâm thỏa (1);
  \item hai thói quen của cùng hoạt động có tâm cách nhau ít nhất $2w$ (cửa sổ không chồng lấn).
\end{enumerate}
Thói quen được mô tả bằng: người $\rho(a)$; hoạt động $a$; giờ điển hình (trung bình vòng tròn các
lần thực hiện trong cửa sổ, mỗi ngày lấy lần gần tâm nhất); khung giờ chứa 80\% số lần; thời lượng
trung vị; số ngày; tần suất $N(c)/|\Dd_\tau|$ và loại ngày $\tau$.
\end{definition}

Định nghĩa giữ đúng tinh thần đề tài thực tập (``lệch không quá khoảng 30 phút'', ``ít nhất 50\%
số ngày'') nhưng thêm chiều \emph{người} và \emph{hoạt động}, và có ba lựa chọn thiết kế quan trọng:

\begin{itemize}
  \item \textbf{Đếm theo ngày, không theo lần.} Một hoạt động lặp năm lần trong một buổi sáng chỉ được
  tính một ngày. Nhờ đó các hoạt động diễn ra nhiều lần trong ngày (vệ sinh, xem TV) không làm phồng
  tần suất.
  \item \textbf{Cửa sổ có bề rộng cố định.} Khác với DBSCAN, các lần thực hiện không thể ``nối
  chuỗi'' qua nhau: khung giờ của một thói quen không thể rộng hơn $2w$. Mục~\ref{sec:exp-ablation}
  cho thấy DBSCAN và HDBSCAN lệch khỏi định nghĩa này đáng kể ngay cả khi dùng nhãn thật.
  \item \textbf{Loại ngày.} Thói quen được tìm trên cả ba loại ngày. Thói quen chỉ thấy ở ngày thường
  hoặc cuối tuần được giữ khi không có thói quen ``tất cả'' nào trong vòng 60 phút và có ít nhất 4
  ngày (tránh trường hợp 3/6 ngày cuối tuần đã đạt 50\%).
\end{itemize}

Thuật toán~\ref{alg:scan} hiện thực định nghĩa bằng cách quét tâm $c$ trên lưới 5 phút.

\begin{algorithm}[htbp]
\caption{Tìm thói quen của một hoạt động trên một loại ngày (quét cửa sổ)}
\label{alg:scan}
\KwIn{giờ bắt đầu $h_1..h_n$ và ngày $d_1..d_n$ của các lần thực hiện hoạt động $a$;
số ngày quan sát $D$; nửa cửa sổ $w$; ngưỡng $\phi$}
\KwOut{danh sách thói quen của $a$}
$G \leftarrow \{0, 5, 10, \dots, 1435\}$ phút \tcp*{lưới tâm}
\ForEach{$c \in G$}{
  $N(c) \leftarrow \big|\{ d_j : |\delta(h_j, c)| \le w \}\big|$ \tcp*{số ngày phân biệt}
}
$Q \leftarrow \{ c \in G : N(c) \ge \phi D\}$\;
Tách $Q$ thành các đoạn liên tục trên vòng tròn; mỗi đoạn cho ứng viên $(c^\ast, N(c^\ast))$ với
$c^\ast$ là điểm giữa vùng đạt $N$ lớn nhất\;
Duyệt ứng viên theo $N$ giảm dần, nhận $c^\ast$ nếu cách mọi tâm đã nhận ít nhất $2w$\;
\ForEach{tâm được nhận $c$}{
  Mỗi ngày lấy lần thực hiện gần $c$ nhất trong $[c-w,c+w]$\;
  Tính giờ điển hình, khung 80\%, thời lượng trung vị, tần suất\;
}
\end{algorithm}

Độ phức tạp là $O(|G| \cdot n)$ với $|G| = 288$, đủ nhanh cho mọi nhà (dưới một giây).

\section{Tổng quan quy trình}
\label{sec:m-overview}

Hình~\ref{fig:pipeline} tóm tắt quy trình. Hàng trên là hướng tiếp cận đề xuất: nhật ký cảm biến được
tiền xử lý, một bộ nhận dạng gán nhãn (người, hoạt động) cho từng sự kiện, các sự kiện được giải mã thành
lần thực hiện, rồi Thuật toán~\ref{alg:scan} rút ra thói quen của từng người. Hàng dưới là các khối dùng
chung và các baseline: cấu trúc nhà rút ra từ dữ liệu (dùng cho cả bộ nhận dạng lẫn baseline), các baseline
không dùng nhãn B0/B1, khâu đánh giá với thói quen đáp án và ứng dụng phát hiện ngày bỏ thuốc. Nhãn thật chỉ
được dùng ở hai chỗ: huấn luyện bộ nhận dạng (ngày huấn luyện) và tạo thói quen đáp án (ngày kiểm tra).

\begin{figure}[htbp]
\centering
\begin{tikzpicture}[
  box/.style={draw, rounded corners=3pt, align=center, font=\footnotesize, minimum height=1.25cm,
              text width=2.35cm, inner sep=3pt, fill=#1},
  box/.default=white,
  arr/.style={-{Stealth[length=2mm]}, thick},
  darr/.style={-{Stealth[length=2mm]}, thick, dashed}]
% Hàng trên: hướng tiếp cận đề xuất
\node[box=gray!10]   (raw)  at (0,0)    {Nhật ký cảm biến\\(luồng trộn\\hai người)};
\node[box=gray!10]   (pre)  at (3.05,0) {Tiền xử lý\\(mục~\ref{sec:m-pre})};
\node[box=blue!10]   (rec)  at (6.1,0)  {Nhận dạng\\(người, hoạt động)\\B2 / M1 / M2};
\node[box=blue!10]   (dec)  at (9.15,0) {Giải mã thành\\lần thực hiện\\(mục~\ref{sec:m-rec})};
\node[box=orange!15] (hab)  at (12.2,0) {Khai phá thói quen\\(Thuật toán~\ref{alg:scan})};
% Hàng dưới: khối dùng chung, baseline, đánh giá
\node[box=gray!10]   (topo) at (3.05,-2.3) {Cấu trúc nhà:\\đồ thị kề, vùng\\(mục~\ref{sec:m-topo})};
\node[box=gray!10]   (base) at (6.1,-2.3)  {Baseline B0 / B1\\phiên theo vùng\\(mục~\ref{sec:m-base})};
\node[box=green!12]  (eval) at (9.15,-2.3) {Đánh giá với\\thói quen đáp án\\(mục~\ref{sec:m-eval})};
\node[box=green!12]  (rq3)  at (12.2,-2.3) {Phát hiện\\ngày bỏ thuốc\\(mục~\ref{sec:m-rq3})};
\draw[arr] (raw) -- (pre);
\draw[arr] (pre) -- (rec);
\draw[arr] (rec) -- (dec);
\draw[arr] (dec) -- (hab);
\draw[arr] (pre) -- (topo);
\draw[arr] (topo) -- (rec);
\draw[darr] (topo) -- (base);
\draw[darr] (base) -- (eval);
\draw[arr] (hab) -- (eval);
\draw[arr] (hab) -- (rq3);
\end{tikzpicture}
\caption{Quy trình của khóa luận. Hàng trên: hướng tiếp cận đề xuất; nét đứt: nhánh baseline không dùng
nhãn. Nhãn thật chỉ dùng để huấn luyện bộ nhận dạng (ngày huấn luyện) và tạo thói quen đáp án (ngày kiểm
tra).}
\label{fig:pipeline}
\end{figure}

\section{Tiền xử lý}
\label{sec:m-pre}

\begin{enumerate}
  \item \textbf{Đọc và sắp xếp.} Một số sự kiện trong CASAS bị ghi lệch thứ tự tối đa 0,1 giây (ghi
  từ nhiều cảm biến cùng lúc); dữ liệu được sắp lại theo thời gian.
  \item \textbf{Ghép nhãn.} Các mốc \texttt{begin}/\texttt{end} cùng tên được ghép thành lần thực
  hiện; một dòng lỗi trong Tulum2 chỉ có chữ \texttt{begin} mà không có tên hoạt động bị bỏ qua.
  \item \textbf{Lọc ngày.} Bỏ ngày có số sự kiện dưới 20\% trung vị số sự kiện mỗi ngày (ngày hỏng
  thiết bị hoặc vắng nhà).
  \item \textbf{Đơn vị quan sát.} Chỉ giữ sự kiện cảm biến chuyển động bật (\texttt{M* ON}): cảm biến
  nhiệt độ đổi chậm và không chỉ ra vị trí người, còn sự kiện \texttt{OFF} chỉ lặp lại thông tin.
  \item \textbf{Chia dữ liệu.} 60\% số ngày đầu để huấn luyện, 40\% số ngày sau để kiểm tra; trong
  tập huấn luyện, 25\% số ngày cuối làm tập kiểm định để chọn ngưỡng và số epoch. Không xáo trộn
  ngẫu nhiên.
  \item \textbf{Tập hoạt động mô hình hóa.} Các hoạt động có ít nhất 5 lần thực hiện trong tập huấn luyện.
\end{enumerate}

\section{Cấu trúc nhà rút ra từ dữ liệu}
\label{sec:m-topo}

Bộ dữ liệu không ghi cảm biến nằm ở phòng nào, và sơ đồ mặt bằng chỉ có dạng ảnh, thậm chí thiếu
tầng hai của Tulum. Vì vậy khóa luận rút cấu trúc nhà từ chính dữ liệu huấn luyện:

\textbf{Đồ thị kề.} Đếm $T_{ij}$ là số lần sự kiện tại $s_i$ được theo sau trong vòng 5 giây bởi sự
kiện tại $s_j \neq s_i$, rồi đối xứng hóa. Hai cảm biến \emph{kề} nhau nếu, theo ít nhất một chiều,
cặp đó chiếm từ 8\% số lần rời khỏi cảm biến. Nếu đồ thị bị tách rời, nối các thành phần bằng cạnh
mạnh nhất. Khoảng cách $\mathrm{hop}(s_i, s_j)$ là số bước ngắn nhất trên đồ thị.

Ngưỡng 8\% và cửa sổ 5 giây được chọn để loại \emph{chuyển tiếp giả}: khi hai người cùng hoạt động,
hai sự kiện liên tiếp của hai người ở hai nơi xa nhau tạo ra cặp ``chuyển tiếp'' không có thật. Với
ngưỡng 2\% và cửa sổ 10 giây, cảm biến bếp M021 của Cairo bị coi là kề cảm biến phòng ngủ tầng trên
M005; với lựa chọn cuối cùng hai cảm biến này cách nhau 5 bước, phù hợp sơ đồ thật.

\textbf{Vùng.} Phân cụm phổ~\cite{vonluxburg2007spectral} trên ma trận tỉ lệ chuyển tiếp đối xứng,
số vùng chọn theo khoảng cách trị riêng trong khoảng 3--10. Vùng đóng vai trò ``phòng'' của phương
pháp thực tập. Trên Cairo, phân vùng tự động khớp với phân vùng đọc từ sơ đồ chính
thức~\cite{casas2025floorplans} ở mức ARI được báo cáo trong Chương~\ref{chap:exp}.

\section{Nhận dạng (người, hoạt động)}
\label{sec:m-rec}

\subsection{Bài toán đa nhãn theo sự kiện}

Với mỗi sự kiện $e_i$, nhãn là vector $y_i \in \{0,1\}^{|\mathcal A|}$ với $y_{ik} = 1$ nếu $t_i$
nằm trong một lần thực hiện hoạt động $a_k$. Vì tên hoạt động mang tiền tố người, ``đúng người'' được
mã hóa sẵn trong nhãn; vì hai người có thể làm hai việc cùng lúc, một sự kiện có thể mang nhiều nhãn.
Mọi bộ nhận dạng chỉ nhìn \emph{quá khứ}: cửa sổ kết thúc tại sự kiện hiện tại, như khi triển khai thật.

\subsection{Ba bộ nhận dạng}

\textbf{B2 --- tra bảng.} Ước lượng $P(a_k \mid \text{vùng của } s_i, \text{khung 30 phút của } t_i)$
bằng tần suất trên dữ liệu huấn luyện. Đây là baseline không học đặc trưng: nó chỉ biết ``ở đâu, lúc
nào'' --- đúng thông tin mà phương pháp thực tập dùng.

\textbf{M1 --- Random Forest.} Với cửa sổ $W = 20$ sự kiện kết thúc tại $e_i$, đặc trưng gồm:
tỉ lệ số lần mỗi cảm biến xuất hiện trong cửa sổ; cảm biến hiện tại (one-hot); vùng hiện tại; $\sin$ và
$\cos$ của giờ; cờ cuối tuần; độ dài cửa sổ (log giây); khoảng cách tới sự kiện trước (log giây); số
cảm biến phân biệt; và \textbf{tỉ lệ ``nhảy xa''} --- số cặp sự kiện liên tiếp cách nhau dưới 10 giây
nhưng ở hai cảm biến cách $\ge 3$ bước trên đồ thị kề. Đặc trưng cuối cùng là dấu hiệu trực tiếp cho
thấy có hơn một người đang hoạt động. Mô hình là rừng 150 cây~\cite{breiman2001rf}, mỗi lá ít nhất 3
mẫu, trọng số lớp cân bằng theo từng mẫu bootstrap, nhiều đầu ra.

\textbf{M2 --- Bi-GRU.} Chuỗi $W = 30$ sự kiện gần nhất; mỗi bước gồm vector nhúng 16 chiều của cảm
biến và ba đặc trưng thời gian ($\sin$, $\cos$ của giờ, $\log(1+\Delta t)/5$). Một lớp GRU hai chiều
64 đơn vị~\cite{cho2014gru,schuster1997brnn}, ghép trạng thái cuối hai chiều, dropout 0,2, tầng tuyến
tính và hàm sigmoid cho từng hoạt động. Hàm mất mát là entropy chéo nhị phân có trọng số lớp dương
$\min(50, \max(1, (1-p_k)/p_k))$ với $p_k$ là tỉ lệ dương của hoạt động $k$. Tối ưu Adam, tốc độ học
$10^{-3}$, lô 512, tối đa 15 epoch, dừng sớm sau 3 epoch không cải thiện mất mát trên tập kiểm định.
Cấu hình được chọn \emph{theo mất mát trên tập kiểm định của Cairo} trong ba cấu hình thử (chi tiết
ở Phụ lục~\ref{app:gru}); sau đó mô hình được huấn luyện lại trên toàn bộ ngày huấn luyện với số epoch
tốt nhất tìm được.

\subsection{Ngưỡng và giải mã}

Xác suất dự đoán được làm trơn bằng trung bình trượt nhân quả trên 5 sự kiện. Với mỗi hoạt động, ngưỡng
quyết định được chọn trong lưới $\{0{,}10; 0{,}15; \dots; 0{,}90\}$ để F1 trên tập kiểm định cao nhất.
Mô hình cuối cùng (B2, M1) được huấn luyện lại trên toàn bộ ngày huấn luyện và dùng ngưỡng đã chọn.

Dự đoán nhị phân theo sự kiện được \emph{giải mã} thành lần thực hiện: các sự kiện dương liên tiếp của
cùng hoạt động, nối khi cách nhau không quá 10 phút, bỏ đoạn có dưới 2 sự kiện. Thời điểm bắt đầu của
lần thực hiện là thời điểm sự kiện đầu tiên. Thói quen dự đoán $\hat\Hb$ thu được bằng cách chạy
Thuật toán~\ref{alg:scan} trên các lần thực hiện đã giải mã.

\section{Các baseline không dùng nhãn}
\label{sec:m-base}

\textbf{B0 --- phương pháp của đề tài thực tập.} Gom sự kiện thành \emph{phiên}: chuỗi sự kiện liên
tiếp ở cùng một vùng, tách khi đổi vùng hoặc im lặng quá 30 phút; bỏ phiên ngắn hơn 1 phút. Mỗi phiên
là một điểm $(R\cos\theta, R\sin\theta, 0{,}5\log(1+d))$ với $R = 24/(2\pi)$ và $d$ là thời lượng
(phút); DBSCAN với $\varepsilon = 0{,}5$ và \texttt{min\_samples} $= \max(5, \lceil 0{,}2 D\rceil)$
chạy riêng từng vùng; cụm lặp lại ở ít nhất 50\% số ngày là thói quen.

\textbf{B1 --- B0 cộng tách luồng theo vị trí.} Lấy ý tưởng của Chen~\cite{chen2021fcmcpd}: hai sự
kiện liên tiếp ở hai nơi xa nhau khó do cùng một người gây ra. Thuật toán~\ref{alg:track} gán mỗi sự
kiện vào một \emph{vệt}; phiên được gom trong từng vệt rồi xử lý như B0.

\begin{algorithm}[htbp]
\caption{Tách luồng sự kiện thành vệt theo vị trí (B1)}
\label{alg:track}
\KwIn{sự kiện $(t_i, s_i)$ theo thời gian; khoảng cách hop; số người $K = 2$; ngưỡng gần $\eta = 2$ bước;
thời gian chờ $\Delta = 15$ phút}
\KwOut{mã vệt của từng sự kiện}
$\mathcal T \leftarrow \emptyset$ \tcp*{mỗi vệt: (mã, cảm biến cuối, thời điểm cuối)}
\ForEach{sự kiện $(t_i, s_i)$}{
  $\mathcal T_a \leftarrow$ các vệt có thời điểm cuối trong vòng $\Delta$ trước $t_i$\;
  $k^\ast \leftarrow \arg\min_{k \in \mathcal T_a} \mathrm{hop}(\text{cảm biến cuối}_k, s_i)$ (hòa: vệt gần đây hơn)\;
  \eIf{$\mathcal T_a = \emptyset$ \textbf{hoặc} ($\mathrm{hop} > \eta$ \textbf{và} $|\mathcal T_a| < K$)}{
    mở vệt mới, gán $e_i$ vào vệt đó\;
  }{
    gán $e_i$ vào $k^\ast$\;
  }
  cập nhật cảm biến cuối và thời điểm cuối của vệt được gán\;
}
\end{algorithm}

\textbf{Gán nhãn cho cụm (chỉ để đánh giá).} B0 và B1 không biết tên hoạt động. Để so với đáp án, mỗi
phiên được gán nhãn của lần thực hiện thật có thời gian chồng lấn lớn nhất; cụm được gán nhãn đa số
trong các phiên có nhãn, và bị coi là ``không nhãn'' nếu dưới 20\% số phiên của nó trùng nhãn nào.
Cách gán này \emph{dùng nhãn của tập kiểm tra} nên có lợi cho baseline; so sánh vì thế thận trọng đối
với các mô hình đề xuất.

\section{Đánh giá thói quen}
\label{sec:m-eval}

\textbf{So khớp.} Thói quen dự đoán và thói quen đáp án \emph{cùng hoạt động} được ghép một--một bằng
thuật toán Hungary~\cite{kuhn1955hungarian} trên chênh lệch giờ điển hình (phút); chỉ nhận cặp lệch
không quá 60 phút. Loại ngày không bắt buộc trùng: thói quen ``uống thuốc 06:57 mọi ngày'' vẫn đúng về
bản chất với đáp án ``uống thuốc 07:08 ngày thường'', vì ngưỡng 50\% đặt trên hai tập ngày khác nhau.

\textbf{Chỉ số.} Với $n_m$ cặp khớp, $|\hat\Hb|$ thói quen dự đoán và $|\Hb^\star|$ thói quen đáp án:
\begin{equation}
  P = \frac{n_m}{|\hat\Hb|}, \qquad R = \frac{n_m}{|\Hb^\star|}, \qquad F1 = \frac{2PR}{P+R}.
\end{equation}
Kèm theo là sai số giờ (trung bình chênh lệch giờ điển hình của các cặp khớp, phút) và sai số tần suất.
Mọi chỉ số được tính hai lần: trên toàn bộ thói quen, và chỉ trên thói quen \textbf{gắn với từng người}
($\rho(a) \in \{R1, R2\}$).

\textbf{Oracle.} Chạy Thuật toán~\ref{alg:scan} trên nhãn thật của ngày kiểm tra cho đúng $\Hb^\star$,
nên F1 của Oracle bằng 1 theo cấu trúc. Điều này có chủ đích: với khâu khai phá cố định là định nghĩa,
\emph{mọi} sai lệch giữa $\hat\Hb$ và $\Hb^\star$ đến từ khâu nhận dạng --- điều kiện cần cho phân
tích lỗi lan truyền.

\section{Phân tích lỗi lan truyền}
\label{sec:m-prop}

Thay vì chỉ quan sát vài mô hình, khóa luận \emph{làm nhiễu nhãn thật có kiểm soát} trên ngày kiểm tra,
khai phá thói quen trên nhãn đã nhiễu và so với đáp án. Bốn loại lỗi mô phỏng những gì một bộ nhận dạng
có thể làm sai:

\begin{itemize}
  \item \textbf{Bỏ sót}: xóa ngẫu nhiên tỉ lệ $p \in \{0; 0{,}1; \dots; 0{,}5\}$ lần thực hiện.
  \item \textbf{Lệch giờ}: dời mỗi lần thực hiện một lượng theo phân phối chuẩn độ lệch
  $\sigma \in \{0, 5, 10, 20, 30, 45\}$ phút.
  \item \textbf{Chèn giả}: thêm số lần thực hiện giả bằng $p \in \{0; 0{,}1; 0{,}25; 0{,}5; 1; 2\}$ lần
  số thật, ở ngày và giờ ngẫu nhiên.
  \item \textbf{Tráo người}: với xác suất $p \in \{0; 0{,}1; \dots; 0{,}5\}$, đổi \texttt{R1} $\leftrightarrow$
  \texttt{R2} trên các hoạt động gắn người --- lỗi đặc trưng của nhà nhiều người; F1 tính trên thói quen
  gắn người.
\end{itemize}
Mỗi mức lặp lại với 5 hạt giống ngẫu nhiên.

\section{Phát hiện bỏ thuốc}
\label{sec:m-rq3}

\begin{enumerate}
  \item \textbf{Thói quen chuẩn} của hoạt động đích (Cairo: \texttt{R2\_Take\_Medicine}; Milan:
  \texttt{Morning\_Meds}, \texttt{Eve\_Meds}) học từ nhãn của các ngày huấn luyện bằng chính
  Định nghĩa~\ref{def:habit}. \emph{Cửa sổ phát hiện} là khung 80\% của thói quen nới thêm 30 phút mỗi phía.
  \item \textbf{Cài bất thường}: chọn ngẫu nhiên 20\% số ngày kiểm tra có uống thuốc, xóa mọi sự kiện cảm
  biến nằm trong các lần uống thuốc của ngày đó --- mô phỏng ``hôm đó không ai tới chỗ để thuốc''. Các sự
  kiện khác, kể cả bữa sáng ở cùng bếp, được giữ nguyên.
  \item \textbf{Phát hiện}: chạy bộ nhận dạng (đã huấn luyện, không học lại) trên luồng đã sửa; nếu
  trong cửa sổ không nhận ra lần thực hiện hoạt động đích nào thì báo động.
  \item \textbf{Baseline theo vùng} (B0-vùng): báo động nếu vùng chứa cảm biến hay kích hoạt nhất khi uống
  thuốc không có sự kiện nào trong cửa sổ.
  \item \textbf{Đáp án} của một ngày là ``bỏ thuốc'' nếu sau khi cài không còn lần uống thuốc thật nào bắt
  đầu trong cửa sổ, gồm cả ngày bị cài và ngày tự nhiên không uống.
\end{enumerate}
Chỉ số là precision, recall, F1 theo ngày, số báo động nhầm trên mỗi ngày và độ trễ báo động (từ giờ điển
hình tới cuối cửa sổ). Kết quả lấy trung bình 5 lần cài với hạt giống khác nhau.
